\scanpage{317}
\mspage{16}
% archive sheet number: 318

(c'est ce qu'on va vérifier directement dans
certains cas remarquables, p.ex. celui où
$L=\varnothing$ et où on trouve la \inserted{ss} catégorie pleine
de $\mathit{Ab}_{\ast}$ engendrée par \inserted{les} $G_0[n]$ grâce aux
seules sommes finies.

c) On a oublié de noter que \deleted{dans}
dans le cas particulier $B=\mathit{Ab}$,
on trouve les $\overline{\mathrm m}$s types d'objets
(savoir ceux de $\mathit{Ab}_{\ast}$ ou de $\mathit{Ab}_{\cdot}$ au
choix) pour exprimer les différents
\uline{inv.} contravariants \inserted{$F$} possibles à valeurs
dans $\mathit{Ab}$ (commutant aux $\varprojlim$) et
les \uline{op.} \inserted{$\Omega$} qu'on peut faire sur des
inv. contravar. \deleted{quelconques} \inserted{quelc.} (à val.
dans une cat \deleted{\unsure{avec}} additive à $\varprojlim$).

On a déj\unsure{à} dit comment
\[
\begin{aligned}
F\ \text{correspondent}&\quad K_{\ast}\in\mathit{Ab}_{\ast}\ \text{ou}\ K_{\cdot}\in\mathit{Ab}_{\cdot}\\
\Omega\ \deleted{\text{\unsure{co}}}\text{\rule{2.2cm}{0.3pt}}&\quad
L_{\ast}\in\mathit{Ab}_{\ast}\ \text{ou}\ L_{\cdot}\in\mathit{Ab}_{\cdot}
\end{aligned}
\]

comment alors exprimer directement la correspondance entre
$F$ et $\Omega$, si $K=L$ ? On trouve
si $F_0\in((\mathit{Ab}_{\ast})^{\circ})_{\ast}=B_{0\ast}$
\inserted{$\simeq(\mathit{Ab}_{\ast}^{\ast})^{\circ}$} est le \deleted{\unsure{constituant|constante}} \inserted{l'invariant}
universel comme dessus, à valeurs dans
$B_0=(\mathit{Ab}_{\ast})^{\circ}$, alors pour tt $X_{\ast}\in\mathit{Ss}=\widehat\Delta$
\[
\boxed{F(X_{\ast})\simeq\Omega_{\deleted{\unsure{\ast}}}^{\ast}(\mathbb C_{\ast}(X_{\ast})^{\circ})\simeq
\Omega_{\unsure{\cdot}}^{\cdot}(\mathbb C_{\cdot}(X_{\ast})^{\circ})}
\]

où l'exposant $\circ$ à $\mathbb C_{\ast}(X_{\ast})$ ou $\mathbb C_{\cdot}(X_{\ast})$
signifie le $\overline{\mathrm m}$, mais considéré comme
objet non de $\mathit{Ab}_{\ast}$, \inserted{resp $\mathit{Ab}_{\cdot}$}, mais de \deleted{$(\mathit{Ab}^{\circ}$}
$B'^{\ast}$ \inserted{resp. $B'^{\cdot}$} où $B'=\mathit{Ab}^{\circ}$, de sorte que $\Omega^{\ast}$
\deleted{resp. $\Omega^{\cdot}$ est} défini sur \deleted{$\mathit{Ab}_{\ast}$ resp. $\unsure{\mathit{A}}$}
$B'^{\ast}$ resp $B'^{\cdot}$ ($\Omega_{B}^{\ast}:B'^{\ast}\longrightarrow B^{\prime}$,
$\Omega^{\cdot}_{B}:B'^{\cdot}\longrightarrow B'$)
a une valeur sur $\mathbb C_{\ast}(X_{\ast})^{\circ}$, \inserted{resp. $\mathbb C_{\cdot}(X_{\ast})$}, et
c'est cet élt de $B=(\mathit{Ab})^{\circ}$ qui est
$F(X_{\ast})$ (à isom. can. près).

Tous ces développements sont encore
valables pour une catégorie $\Delta$ (ess. petite)
quelconque -- et aussi dans le contexte
non additif -- \deleted{la seule chose} \inserted{Sauf qu'on} met des
$\ast$ et non un $\cdot$, qui, eux, nécessitent
des catégories de valeurs additives et le
\unsure{topos|topo} bien spécial de Dold-Puppe.

En fait (revenant au régime $L$
de cat. d'indices envisagé dans b))
on s'intéressera surtout aux
inv. contravariants en $X_{\ast}$ provenant
d'objets $K_{\ast}\in\mathit{Ab}_{\ast}$ qui sont dans $U_L$ --
i.e. qui peuvent se déduire des
complexes $C_{\ast}^{[n]}\in(\mathit{Ab})_{\ast}$ (et des opérations
$\tfrac12$ simpliciales entre lesdites) à l'aide des
opérations $\Omega_{\ast}$ \deleted{\unsure{associées} \unsure{à des}} \inserted{provenant par $\varprojlim$, soit des $c^{n}_{\ast}\in(\mathit{Ab})_{\ast}$,}
et des op. $\tfrac12$ simpl. entre \unsure{icix|iceux}, via des
\deleted{telles que} $\varprojlim$ des $\overline{\mathrm m}$ de
dire que \deleted{des} on
s'intéressera aux
foncteurs \deleted{\unsure{$F(X_{\ast})$}}

% slanted
\margin{On peut bien entendu, comme d'habitude, introduire un $k$, et remplacer partout $\mathit{Ab}$ par $\mathit{Ab}_{k}$ --}
\margin{($\mathit{Ab}$) -- Attention, dans cette formule apparaît \unsure{le caractère en autodualité|la confusion en attendant} de la situation -- il s'agit bien des $F$ \uline{contravariants} à val. dans $\mathit{Ab}$ lui-$\overline{\mathrm m}$ (et non dans $(\mathit{Ab})^{\circ}$ --)}
